\scanpage{63}
On \deleted{\unsure{voit|veut}} se demande s'il est représentable, i.e. on demande la valeur en $T$ du \uline{foncteur adjoint} (au sens de Kan) du foncteur $\underline{\mathbb W}$ des $S$-préschémas en $T$-préschémas. Une solution du Pb universel posé par le foncteur en $T$ ci-dessus est par définition un couple
\[
 X\quad(S\text{-préschéma})\qquad \underline{\mathbb W}(X)\xrightarrow{\alpha}\unsure{\mathfrak X|\mathcal X}
\]
tel que pour tout $T$, l'application naturelle
\[
 \mathrm{Hom}_S(T,X)\xrightarrow{\sim}\mathrm{Hom}_{\mathbb W(S)}(\mathbb W(T),\mathfrak X)
\]
\deleted{soit une} \unsure{définie|obtenue} par factorisation
\[
 \underline{\mathbb W}(T)\xrightarrow{\mathbb W(u)}\underline{\mathbb W}(X)\xrightarrow{\alpha}\mathfrak X
\]
soit \uline{bijective}. Comme le foncteur $(*)$ en $T$ est de ``nature locale'', il suffit d'ailleurs de vérifier la bijectivité pour $T$ affine \deleted{\unsure{et}} \inserted{sur un ouvert} affine \deleted{\unsure{de|sur}} $S$, p. ex.

Désignons par $\underline{\mathbb W}^a$ le foncteur adjoint de $\underline{\mathbb W}$ (défini sur une ss-catégorie de la catégorie des $T$-préschémas)
\[
 \mathrm{Hom}_{S^a}(T^a,\mathfrak X)\simeq\mathrm{Hom}_S(T,\mathfrak X
\]
% archive sheet number: 62
